\documentclass[10pt]{article}

% Arxiv style format


% \usepackage{tgpagella}
\usepackage{parskip}


% \renewcommand{\paragraph}[1]{
% 	\textbf{#1}~~}

\usepackage[letterpaper, left=1in, right=1in, top=1in, bottom=1in]{geometry}

\usepackage[dvipsnames]{xcolor}
\usepackage[colorlinks=true, linkcolor=blue!70!black, citecolor=blue!70!black]{hyperref}
\usepackage{microtype}


% \usepackage{natbib}
% \bibliographystyle{alpha}
% \bibpunct{(}{)}{;}{a}{,}{,}

\usepackage{amsthm}
\usepackage{mathtools}
\usepackage{amsmath}
\usepackage{bm}
\usepackage{bbm}
\usepackage{amsfonts}
\usepackage{amssymb}
\let\vec\undefined
% \usepackage{MnSymbol} %Actually conflicts with amssymb and others
% yair: I don't like how MnSymbol looks like. If others love it we can put it
% back for the arxiv version, but for COLT perhaps we should avoid it to be
% consistent with the PMLR formatting?
\usepackage{xpatch}
\usepackage{pifont}
\usepackage{array}
\usepackage{booktabs}
\usepackage{floatrow}
% Table float box with bottom caption, box width adjusted to content
\newfloatcommand{capbtabbox}{table}[][\FBwidth]
\usepackage{caption}
\usepackage{subcaption}

\usepackage{algorithm}
\usepackage{algorithmicx}
\usepackage{algpseudocode}

\usepackage{comment}

% \usepackage{enumitem}
% \usepackage{amsmath}
% \usepackage{amsthm}
\usepackage{bbm}

% Align Block
\newcommand\numberthis{\addtocounter{equation}{1}\tag{\theequation}}
\allowdisplaybreaks

\usepackage{mathtools}

% Math delimiters
\DeclarePairedDelimiter{\abs}{\lvert}{\rvert}
\DeclarePairedDelimiter{\brk}{[}{]}
\DeclarePairedDelimiter{\crl}{\{}{\}}
\DeclarePairedDelimiter{\prn}{(}{)}
\DeclarePairedDelimiter{\nrm}{\|}{\|}
\DeclarePairedDelimiter{\tri}{\langle}{\rangle}
\DeclarePairedDelimiter{\dtri}{\llangle}{\rrangle}
\DeclarePairedDelimiter{\ceil}{\lceil}{\rceil}
\DeclarePairedDelimiter{\floor}{\lfloor}{\rfloor}
\newcommand{\midsem}{\,;}
\newcommand{\trc}[1]{\text{tr} \prn*{#1}}

\let\Pr\undefined
\let\P\undefined
\DeclareMathOperator{\En}{\mathbb{E}}
\newcommand{\Ep}{\wt{\bbE}}
%\DeclareMathOperator*{\Eh}{\widehat{\mathbb{E}}}
\DeclareMathOperator{\P}{P}
\DeclareMathOperator{\Pr}{\bbP}

% Arg<x>
\DeclareMathOperator*{\argmin}{argmin} % * Places subscript directly under operator
\DeclareMathOperator*{\argmax}{argmax}
\DeclareMathOperator*{\arginf}{arginf}
\DeclareMathOperator*{\argsup}{argsup}

% one-off macros
\newcommand{\ls}{\ell}
\newcommand{\bls}{\mb{\ls}}
\newcommand{\ind}[1]{\mathbbm{1}\crl*{#1}}    %Indicator
\newcommand{\indd}[1]{\mathbbm{1}\crl{#1}}    %Indicator
\newcommand{\eps}{\varepsilon}
\newcommand{\veps}{\varepsilon}

\newcommand{\ldef}{\vcentcolon=}
\newcommand{\rdef}{=\vcentcolon}

% styles
\newcommand{\mc}[1]{\mathcal{#1}}
\newcommand{\wt}[1]{\widetilde{#1}}
\newcommand{\wh}[1]{\widehat{#1}}
%\newcommand{\mb}[1]{\boldsymbol{#1}}


% Special letters: blackboard, mathcal, widehat
% djhsu magic
\def\ddefloop#1{\ifx\ddefloop#1\else\ddef{#1}\expandafter\ddefloop\fi}
\def\ddef#1{\expandafter\def\csname bb#1\endcsname{\ensuremath{\mathbb{#1}}}}
\ddefloop ABCDEFGHIJKLMNOPQRSTUVWXYZ\ddefloop
\def\ddefloop#1{\ifx\ddefloop#1\else\ddef{#1}\expandafter\ddefloop\fi}
\def\ddef#1{\expandafter\def\csname b#1\endcsname{\ensuremath{\mathbf{#1}}}}
\ddefloop ABCDEFGHIJKLMNOPQRSTUVWXYZ\ddefloop
\def\ddef#1{\expandafter\def\csname c#1\endcsname{\ensuremath{\mathcal{#1}}}}
\ddefloop ABCDEFGHIJKLMNOPQRSTUVWXYZ\ddefloop
\def\ddef#1{\expandafter\def\csname h#1\endcsname{\ensuremath{\widehat{#1}}}}
\ddefloop ABCDEFGHIJKLMNOPQRSTUVWXYZabcdefghijklmnopqrsuvwxyz\ddefloop    % Not defined for t.
\def\ddef#1{\expandafter\def\csname hc#1\endcsname{\ensuremath{\widehat{\mathcal{#1}}}}}
\ddefloop ABCDEFGHIJKLMNOPQRSTUVWXYZ\ddefloop
\def\ddef#1{\expandafter\def\csname t#1\endcsname{\ensuremath{\widetilde{#1}}}}
\ddefloop ABCDEFGHIJKLMNOPQRSTUVWXYZ\ddefloop
\def\ddef#1{\expandafter\def\csname tc#1\endcsname{\ensuremath{\widetilde{\mathcal{#1}}}}}
\ddefloop ABCDEFGHIJKLMNOPQRSTUVWXYZ\ddefloop

% Names
\newcommand{\Holder}{H{\"o}lder}

% COMMONLY USED MACROS
%Losses
\newcommand{\KL}[2]{\mathrm{KL}{\prn*{#1 \| #2}}}
\newcommand{\kl}[2]{\mathrm{kl}{\prn*{#1 \| #2}}}
\newcommand{\ball}{\mathbb{B}}
\newcommand{\poprisk}{L_{\cD}}
\newcommand{\emprisk}{\wh{L}_n}
\newcommand{\sign}{\mathrm{sign}}

% Matrices
%\newcommand{\diag}{\textrm{diag}}
%\newcommand{\rank}{\mathrm{rank}}
\newcommand{\trn}{\dagger}
\newcommand{\Tr}{\mathbf{Tr}}

% Special Symbols
\newcommand{\xr}[1][n]{x_{1:#1}}
\newcommand{\br}[1][n]{b_{1:#1}}
\newcommand{\er}[1][n]{\eps_{1:#1}}
\newcommand{\sr}[1][n]{\sigma_{1:#1}}
\newcommand{\yr}[1][n]{y_{1:#1}}
\newcommand{\zr}[1][n]{z_{1:#1}}
\renewcommand{\th}{{\text{th}}}
%Circled Numbers
\usepackage{tikz}
\newcommand*\circled[1]{\tikz[baseline=(char.base)]{
            \node[shape=circle,draw,inner sep=2pt] (char) {#1};}}

% Calculus
\newcommand{\grad}{\nabla}

\newcommand{\proman}[1]{\prn*{\romannumeral #1}}
\newcommand{\overles}[1]{\overset{ #1}{\lesssim{}}}
\newcommand{\overleq}[1]{\overset{ #1}{\leq{}}}
\newcommand{\overgeq}[1]{\overset{#1}{\geq{}}}
\newcommand{\overeq}[1]{\overset{#1}{=}}
\newcommand{\overapx}[1]{\overset{#1}{\approx}}

%%% theorems

%\theoremstyle{definition}  %Sets style of subsequent newtheorems to 'definition'
%\newtheorem{assumption}{Assumption}
\newtheorem{condition}{Condition}
\newtheorem{claim}{Claim}
%\newtheorem{exercise}{Exercise}
\newtheorem{assumption}{Assumption}
\newtheorem{conjecture}{Conjecture}
\newtheorem{corollary}{Corollary}
\newtheorem{proposition}{Proposition}
%\newtheorem{fact}{Fact}
%\newtheorem{goal}{Goal}
%%\newtheorem{model}{Model}
\newtheorem{problem}{Problem}
%\newtheorem{model}{Model}
%\theoremstyle{plain}
{\newtheorem{lemma}{Lemma}}
\newtheorem{remark}{Remark}
%\newtheorem{example}{Example}
{\newtheorem{theorem}{Theorem}}
\newtheorem{theorem*}{Theorem}
\newtheorem{definition}{Definition}
%\newtheorem{question}{Question}
% \newtheorem{proof}{Proof}



% \newtheorem*{assumption*}{\assumptionnumber}
% \providecommand{\assumptionnumber}{}
% \makeatletter
% \newenvironment{spassumption}[1]
%  {%
%   \renewcommand{\assumptionnumber}{Assumption #1}%
%   \begin{assumption*}%
%   \protected@edef\@currentlabel{#1}%
%  }
%  {%
%   \end{assumption*}
%  }
% \makeatother

% \usepackage{xpatch}
% \xpatchcmd{\proof}{\itshape}{\normalfont\proofnameformat}{}{}
% \newcommand{\proofnameformat}{\bfseries}

% The template's prettyref helpers are unused by this method note.
% They are omitted from this copy because prettyref.sty is not available
% in the installed TinyTeX distribution.

\newcommand{\algcomment}[1]{\textcolor{blue!70!black}{\footnotesize{\texttt{\textbf{//
          #1}}}}}
\newcommand{\algcommentbig}[1]{\textcolor{blue!70!black}{\footnotesize{\texttt{\textbf{/*
          #1~*/}}}}}

\newcommand{\var}{\text{Var}}

\newcommand{\HI}{\text{Household income}}

\newcommand{\dif}{\text{d}}

\newcommand{\cov}{\text{Cov}}

\newcommand{\corr}{\text{Corr}}

\title{Two-Step Gated GRPO for Long-Horizon Trajectories}
\author{\vspace{-2em}}
\date{September 26, 2026}
\setlength{\parskip}{4pt}
\AtBeginDocument{%
  \setlength{\abovedisplayskip}{6pt}%
  \setlength{\belowdisplayskip}{6pt}%
  \setlength{\abovedisplayshortskip}{3pt}%
  \setlength{\belowdisplayshortskip}{3pt}%
}

\begin{document}
\maketitle

\begin{abstract}
We specify a two-step gated GRPO method for long-horizon trajectories.
Terminal rewards are averaged within each observed action, distinct action
means are standardized equally, and a near-zero first action receives an
additional signal only when its successor passes a dispersion gate. We define
the resulting advantages and clipped objective; no experiments are reported.
\end{abstract}

\section{Problem and data}

Let \(\pi_{\mathrm{old}}\) be the known behavior policy, \(\tau_i\) an old
trajectory, and \(R_i\) its terminal reward. A state \(s\) denotes the
complete decision prefix; \(a_i(s)\) is the action taken by trajectory \(i\)
at that state. The aim is to identify first actions with near-zero
action-mean advantage whose successor actions show a clearer signed
separation.

\section{Action-mean advantage at a state}

For any observed state \(s\), define
\begin{equation}
\mathcal I(s)=\{i:\tau_i\text{ visits }s\},\qquad
K_s=|\mathcal I(s)|,\qquad
\mathcal U(s)=\{a_i(s):i\in\mathcal I(s)\}.
\end{equation}
For \(u\in\mathcal U(s)\), let
\begin{equation}
\mathcal I(s,u)=\{i\in\mathcal I(s):a_i(s)=u\},\qquad
n_s(u)=|\mathcal I(s,u)|,\qquad
q_s(u)=\frac{1}{n_s(u)}\sum_{i\in\mathcal I(s,u)}R_i.
\end{equation}
Different observed actions then have equal weight, irrespective of their
trajectory counts:
\begin{equation}
d_s=|\mathcal U(s)|,\qquad
\bar q_s=\frac{1}{d_s}\sum_{u\in\mathcal U(s)}q_s(u),\qquad
\sigma_{q,s}=
\left[\frac{1}{d_s}\sum_{u\in\mathcal U(s)}
       \bigl(q_s(u)-\bar q_s\bigr)^2\right]^{1/2}.
\end{equation}
The action-mean advantage is returned to every trajectory with that action:
\begin{equation}
A_i^{\mathrm{state}}(s)=
\begin{cases}
\dfrac{q_s(a_i(s))-\bar q_s}{\sigma_{q,s}},
    &\sigma_{q,s}>0,\\[6pt]
0, &\sigma_{q,s}=0.
\end{cases}
\end{equation}
The zero-variance rule applies to the advantage at this position, not to
the entire trajectory. Repeated actions receive no extra inverse-count
weight in the training objective.

\section{Two selection rounds and the dispersion gate}

At a parent state \(s_t\), use the same action-mean standardization for
the first-round advantage and its dispersion:
\begin{equation}
\sigma_t=\sigma_{q,s_t},\qquad
A_{1,i}(s_t)=A_i^{\mathrm{state}}(s_t),\qquad
M_t=\lceil0.2K_t\rceil.
\end{equation}
If \(\sigma_t=0\), the state creates no gated candidate. Otherwise, let
\(\mathcal L_t(s_t)\) contain exactly \(M_t\) old trajectories with the
smallest \(|A_{1,i}(s_t)|\). A cutoff tie is filled uniformly without
replacement, so only the selected trajectories retain first-step
eligibility. Candidate actions are deduplicated after this trajectory
selection:
\begin{equation}
\mathcal C_t(s_t)=\{a_{i,t}:i\in\mathcal L_t(s_t)\}.
\end{equation}

For each candidate \(a\), set \(s_{t+1}=(s_t,a)\). Use \emph{all} old
trajectories reaching this child state, including those not selected in
the first round. Require at least five distinct observed next actions:
\begin{equation}
d=|\mathcal U(s_{t+1})|\geq5,\qquad
\sigma_{t+1}=\sigma_{q,s_{t+1}}>0.
\end{equation}
The parent and child dispersions are both equally weighted standard
deviations of distinct action-mean rewards. The dispersion gate is
\begin{equation}
G_{\Delta}(s_t,a)
=\mathbf{1}\{\sigma_{t+1}\geq\omega\sigma_t\},
\end{equation}
where \(\omega\) is an unspecified tunable constant. When the gate
passes, compute the signed second-step advantage
\begin{equation}
A_2(s_{t+1},u)
=\frac{q_{s_{t+1}}(u)-\bar q_{s_{t+1}}}{\sigma_{t+1}}.
\end{equation}
Let \(\mathcal S(s_t,a)\) contain the \(m(d)=\lceil d/2\rceil\)
distinct actions with largest \(|A_2(s_{t+1},u)|\). Fill cutoff ties
uniformly at random. If any child condition or the dispersion gate
fails, define \(\mathcal S(s_t,a)=\varnothing\).

\section{Gates and two-step advantages}

The action-level gate and trajectory-level first-step gate are
\begin{align}
G_t^{\mathrm{act}}(s_t,a)
&=\mathbf{1}\{\sigma_t>0\}
  \mathbf{1}\{a\in\mathcal C_t(s_t)\}
  \mathbf{1}\{\mathcal S(s_t,a)\neq\varnothing\},\\
B_{i,t}
&=\mathbf{1}\{i\in\mathcal L_t(s_t)\}
  G_t^{\mathrm{act}}(s_t,a_{i,t}).
\end{align}
Thus \(B_{i,t}=1\) implies that at least one successor action passed
the second round. A same-\(a\) trajectory outside \(\mathcal L_t\)
does not receive the first-step increment
\begin{equation}
\Delta A^{(1)}_{i,t}=\lambda B_{i,t},\qquad\lambda>0.
\end{equation}

For a passing child state and a non-overlapping second-step position,
define \(I^{\mathcal S}_{i,t+1}
=\mathbf{1}\{a_{i,t+1}\in\mathcal S(s_t,a_{i,t})\}\). The second-step
coefficient is
\begin{equation}
\widetilde A_{i,t+1}
=\bigl[1+(\beta-1)I^{\mathcal S}_{i,t+1}\bigr]
 A_2(s_{t+1},a_{i,t+1}),\qquad\beta>1.
\end{equation}
Since \(A_i^{\mathrm{state}}(s_{t+1})=A_2(s_{t+1},a_{i,t+1})\)
at this child state, unselected actions retain their ordinary
action-mean advantage. Selected positive and negative actions both
receive a larger absolute coefficient.

If one position is the previous chain's second action and the next
chain's selected first action, the confirmed overlap rule adds their
two roles:
\begin{equation}
\widehat A_{i,j}^{\mathrm{overlap}}
=A_i^{\mathrm{state}}(s_{i,j})+\lambda
 +\beta I^{\mathcal S}_{i,j}
 A_2(s_{i,j},a_{i,j}).
\end{equation}
When \(I^{\mathcal S}_{i,j}=1\), the action-level part is
\((1+\beta)A_2\), because the two roles are added.

\section{Per-position advantage and GRPO update}

At each trajectory position, let
\begin{equation}
P_{i,j}
=G_{j-1}^{\mathrm{act}}(s_{i,j-1},a_{i,j-1}),\qquad
I^{\mathcal S}_{i,j}
=\mathbf{1}\{a_{i,j}\in
 \mathcal S(s_{i,j-1},a_{i,j-1})\}.
\end{equation}
Take a nonexistent previous gate at a trajectory boundary to be zero.
Abbreviate \(A^{\mathrm{state}}_{i,j}
=A_i^{\mathrm{state}}(s_{i,j})\) and
\(A_{2,i,j}=A_2(s_{i,j},a_{i,j})\) when \(P_{i,j}=1\).
The complete training advantage is
\begin{equation}
\widehat A_{i,j}=
\begin{cases}
A^{\mathrm{state}}_{i,j},
 &P_{i,j}=0,\ B_{i,j}=0,\\
A^{\mathrm{state}}_{i,j}+\lambda,
 &P_{i,j}=0,\ B_{i,j}=1,\\
\beta A_{2,i,j},
 &P_{i,j}=1,\ B_{i,j}=0,\ I^{\mathcal S}_{i,j}=1,\\
A^{\mathrm{state}}_{i,j},
 &P_{i,j}=1,\ B_{i,j}=0,\ I^{\mathcal S}_{i,j}=0,\\
\widehat A_{i,j}^{\mathrm{overlap}},
 &P_{i,j}=1,\ B_{i,j}=1.
\end{cases}
\end{equation}

For the known behavior policy \(\pi_{\mathrm{old}}\), define
\begin{align}
r_{i,j}(\theta)
&=\frac{\pi_\theta(a_{i,j}\mid s_{i,j})}
        {\pi_{\mathrm{old}}(a_{i,j}\mid s_{i,j})},\\
\ell_{\mathrm{clip}}(r,A)
&=\min\{rA,\operatorname{clip}(r,1-\epsilon,1+\epsilon)A\}.
\end{align}
Following the per-trajectory and per-position averaging of
GRPO~\cite{shao2024deepseekmath}, maximize
\begin{equation}
\begin{aligned}
J(\theta)
=\mathbb E_{\{\tau_i\}\sim\pi_{\mathrm{old}}}\Bigg[
 \frac{1}{N}\sum_{i=1}^{N}\frac{1}{T_i}
 \sum_{j=1}^{T_i}\Big(
 &\ell_{\mathrm{clip}}
   (r_{i,j}(\theta),\widehat A_{i,j})\\[-2pt]
 &-\kappa D^{\mathrm{KL}}_{i,j}(\theta)
 \Big)\Bigg].
\end{aligned}
\end{equation}
Equivalently, minimize \(-J(\theta)\). Old rewards, gates, and
advantages are fixed during an update; gradients pass through the
current-policy ratio and an optional KL term. A repeated action
contributes one term per old trajectory occurrence, with no
additional inverse-count correction.

If KL regularization is enabled, the sampled-action estimator used
in~\cite{shao2024deepseekmath} can be written
\begin{equation}
D^{\mathrm{KL}}_{i,j}(\theta)
=\frac{\pi_{\mathrm{ref}}(a_{i,j}\mid s_{i,j})}
       {\pi_\theta(a_{i,j}\mid s_{i,j})}
 -\log\frac{\pi_{\mathrm{ref}}(a_{i,j}\mid s_{i,j})}
            {\pi_\theta(a_{i,j}\mid s_{i,j})}-1.
\end{equation}
The clip width \(\epsilon\), optional KL weight \(\kappa\), and
second-step amplifier \(\beta\) are distinct parameters. Their
numerical values are not set here.

\section{Calculation order}

Collect old trajectories and terminal rewards from the known
\(\pi_{\mathrm{old}}\). Group them by complete state and next action,
then compute \(q_s(u)\), \(\sigma_{q,s}\), and
\(A_i^{\mathrm{state}}(s)\). At each parent state select the lowest
\(20\%\) of trajectories by \(|A_1|\). For each resulting action,
use all old trajectories at its child state; check \(d\geq5\),
nonzero parent and child dispersions, and the \(\Delta\) gate.
Select \(m(d)\) next actions by \(|A_2|\), assign
\(\widehat A_{i,j}\) using the five cases above, and update
\(\theta\) with the clipped objective.

This document specifies a theoretical procedure, not an experiment.
The numerical values of \(\omega\), \(\lambda\), \(\beta\), and the
mapping from an action to a token or full step remain unspecified.

\bibliographystyle{plain}
\bibliography{refs}

\end{document}
